\documentclass[11pt]{article}
\pagestyle{plain}

% \pgfplotsset{compat=1.18}
\usepackage[margin=1in]{geometry}
\usepackage{amsmath,amsfonts,amssymb,amsthm}
\usepackage[ruled,vlined,linesnumbered,commentsnumbered]{algorithm2e}
\SetKwInput{KwIn}{Input}
\SetKwInput{KwOut}{Output}
\SetKwInput{KwGlobal}{Global variable}
\SetKwInput{KwOracle}{Oracle}
\SetKwComment{Comment}{// }{}
\usepackage{graphicx}
\usepackage{enumerate}
\usepackage{bbm}
\usepackage{verbatim}
\usepackage{hyperref,color}
\usepackage[capitalize,nameinlink]{cleveref}
\usepackage[dvipsnames]{xcolor}
\hypersetup{
	colorlinks=true,
	pdfpagemode=UseNone,
	citecolor=OliveGreen,
	linkcolor=NavyBlue,
	urlcolor=Magenta,
	pdfstartview=FitW
}
\usepackage{appendix}
\usepackage{makecell}
\usepackage{tablefootnote}
\crefname{appsec}{Appendix}{Appendices}
\usepackage{tikz}
\usepackage{pgfplots}
\usepackage{xifthen}
\usepackage{subcaption}
\usepackage{hhline}

\theoremstyle{plain}
\newtheorem{theorem}{Theorem}[section]
\newtheorem{proposition}[theorem]{Proposition}
\newtheorem{lemma}[theorem]{Lemma}
\newtheorem{corollary}[theorem]{Corollary}
\newtheorem{conjecture}[theorem]{Conjecture}
\newtheorem{problem}[theorem]{Problem}
\newtheorem{claim}[theorem]{Claim}
\newtheorem{fact}[theorem]{Fact}

\theoremstyle{definition}
\newtheorem{definition}[theorem]{Definition}
\newtheorem{example}[theorem]{Example}
\newtheorem{observation}[theorem]{Observation}
\newtheorem{assumption}[theorem]{Assumption}
\newtheorem*{assumption*}{Assumption}
\newtheorem{condition}[theorem]{Condition}

\theoremstyle{remark}
\newtheorem{remark}[theorem]{Remark}

\crefname{lemma}{Lemma}{Lemmas}
\crefname{theorem}{Theorem}{Theorems}
\crefname{definition}{Definition}{Definitions}
\crefname{fact}{Fact}{Facts}
\crefname{claim}{Claim}{Claims}
\crefname{proposition}{Proposition}{Propositions}


\newcommand{\dif}{\,\mathrm{d}}
%\newcommand{\E}{\mathbb{E}}
%\newcommand{\Var}{\mathrm{Var}}
\newcommand{\Cov}{\mathrm{Cov}}
\newcommand{\MI}{\mathrm{I}}
%\newcommand{\Ent}{\mathrm{Ent}}
\newcommand{\ind}{\mathbbm{1}}
\newcommand{\allone}{\mathbf{1}}
\newcommand{\tr}{\mathrm{tr}}
\newcommand{\wt}{\mathrm{weight}}
\DeclareMathOperator*{\argmax}{arg\,max}
\newcommand{\norm}[1]{\left\lVert #1 \right\rVert}
\newcommand{\TV}[2]{\left\|#1 - #2\right\|_{\mathrm{TV}}}
\newcommand{\ceil}[1]{\left\lceil #1 \right\rceil}
\newcommand{\floor}[1]{\left\lfloor #1 \right\rfloor}
\newcommand{\ip}[2]{\left\langle #1 , #2 \right\rangle}
\newcommand{\grad}{\nabla}
\newcommand{\hessian}{\nabla^2}
\newcommand{\trans}{\intercal}
\newcommand{\diag}{\mathrm{diag}}
\newcommand{\poly}{\mathrm{poly}}
\newcommand{\dist}{\mathrm{dist}}
\newcommand{\sgn}{\mathrm{sgn}}
\newcommand{\fpras}{\mathsf{FPRAS}}
\newcommand{\fpaus}{\mathsf{FPAUS}}
\newcommand{\fptas}{\mathsf{FPTAS}}

\newcommand{\pll}{\parallel}
\newcommand{\eps}{\varepsilon}

\newcommand{\N}{\mathbb{N}}
\newcommand{\Z}{\mathbb{Z}}
\newcommand{\Q}{\mathbb{Q}}
\newcommand{\R}{\mathbb{R}}
\newcommand{\C}{\mathbb{C}}

\renewcommand{\AA}{\mathcal{A}}
\newcommand{\BB}{\mathcal{B}}
\newcommand{\CC}{\mathcal{C}}
\newcommand{\DD}{\mathcal{D}}
\newcommand{\EE}{\mathcal{E}}
\newcommand{\FF}{\mathcal{F}}
\newcommand{\GG}{\mathcal{G}}
\newcommand{\HH}{\mathcal{H}}
\newcommand{\II}{\mathcal{I}}
\newcommand{\JJ}{\mathcal{J}}
\newcommand{\KK}{\mathcal{K}}
\newcommand{\LL}{\mathcal{L}}
\newcommand{\MM}{\mathcal{M}}
\newcommand{\NN}{\mathcal{N}}
\newcommand{\OO}{\mathcal{O}}
\newcommand{\PP}{\mathcal{P}}
\newcommand{\QQ}{\mathcal{Q}}
\newcommand{\RR}{\mathcal{R}}
\renewcommand{\SS}{\mathcal{S}}
\newcommand{\TT}{\mathcal{T}}
\newcommand{\UU}{\mathcal{U}}
\newcommand{\VV}{\mathcal{V}}
\newcommand{\WW}{\mathcal{W}}
\newcommand{\XX}{\mathcal{X}}
\newcommand{\YY}{\mathcal{Y}}
\newcommand{\ZZ}{\mathcal{Z}}

\newcommand{\KL}[2]{D_{\mathrm{KL}} \left( #1 \,\Vert\, #2 \right)}
\newcommand{\KLbig}[2]{D_{\mathrm{KL}} \left( #1 \,\big\Vert\, #2 \right)}
\newcommand{\KLBig}[2]{D_{\mathrm{KL}} \left( #1 \,\Big\Vert\, #2 \right)}

\newcommand{\SKL}[2]{D_{\mathrm{SKL}} \left( #1 \,\Vert\, #2 \right)}
\newcommand{\SKLbig}[2]{D_{\mathrm{SKL}} \left( #1 \,\big\Vert\, #2 \right)}
\newcommand{\SKLBig}[2]{D_{\mathrm{SKL}} \left( #1 \,\Big\Vert\, #2 \right)}

\newcommand{\qandq}{\quad\text{and}\quad}
\newcommand{\supp}{\mathrm{supp}}

\newcommand{\DKL}[2]{\-D_{\-{KL}}\left(#1 \parallel #2\right)}
\newcommand{\DTV}[2]{\-D_{\mathrm{TV}}\left({#1},{#2}\right)}
% \newcommand{\dist}{\mathrm{dist}}
\newcommand{\e}{\mathrm{e}}
\renewcommand{\epsilon}{\varepsilon}
\renewcommand{\emptyset}{\varnothing}
% \newcommand{\norm}[1]{\left\Vert#1\right\Vert}
\newcommand{\set}[1]{\left\{#1\right\}}
\newcommand{\tuple}[1]{\left(#1\right)} 
% \newcommand{\eps}{\varepsilon}
\newcommand{\inner}[2]{\left\langle #1,#2\right\rangle}
\newcommand{\tp}{\tuple}
\newcommand{\ol}{\overline}
\newcommand{\abs}[1]{\left\vert#1\right\vert}
\newcommand{\ctp}[1]{\left\lceil#1\right\rceil}
\newcommand{\ftp}[1]{\left\lfloor#1\right\rfloor}
\newcommand{\Ind}[1]{\-{Ind}_{#1}}
\newcommand{\todo}[1]{{\color{red} TODO: #1}}

\newcommand{\Cor}{\Psi^{\mathrm{sym}}}
\newcommand{\cor}{\Psi^{\-{cor}}}
\newcommand{\Inf}{\Psi}
% \newcommand{\diag}{\-{diag}}

\def\*#1{\boldsymbol{#1}} % Use \*A for \mathbf{A}
\def\+#1{\mathcal{#1}} % Use \+A for \mathcal{A}
\def\-#1{\mathrm{#1}} % Use \-A for \mathrm{A}
\def\=#1{\mathbb{#1}} % Use \^A for \mathbb{A}
\def\!#1{\mathfrak{#1}} % Use \!A for \mathfrak{A}

\newcommand*\diff{\mathop{}\!\mathrm{d}}

% probability
% \DeclareMathOperator*{\oPr}{\mathbf{Pr}}
\def\oPr{\mathop{\mathrm{Pr}}}
\renewcommand{\Pr}[2][]{ \ifthenelse{\isempty{#1}}
  {\oPr\left[#2\right]}
  {\oPr_{#1}\left[#2\right]} } % Use \Pr[a]{b} for \mathbf{Pr}_a[b], \Pr{b} for  \mathbf{Pr}[b]

% expectation: relies on the package "dsfont"
% \DeclareMathOperator*{\oE}{\mathbb{E}}
\def\oE{\mathop{\mathbb{E}}}
\newcommand{\E}[2][]{ \ifthenelse{\isempty{#1}}
  {\oE\left[#2\right]}
  {\oE_{#1}\left[#2\right]} }

% variance
%\DeclareMathOperator*{\oVar}{\mathbf{Var}}
\def\oVar{\mathrm{Var}}
\newcommand{\Var}[2][]{ \ifthenelse{\isempty{#1}}
  {\oVar\left[#2\right]}
  {\oVar_{#1}\left[#2\right]} }

% entropy
% \DeclareMathOperator*{\oEnt}{\mathbf{Ent}}
\def\oEnt{\mathrm{Ent}}
\newcommand{\Ent}[2][]{ \ifthenelse{\isempty{#1}}
  {\oEnt\left[#2\right]}
  {\oEnt_{#1}\left[#2\right]} }

\newcommand{\PhiEnt}[2][]{ \ifthenelse{\isempty{#1}}
  {\oEnt^\phi\left[#2\right]}
  {\oEnt^\phi_{#1}\left[#2\right]} }


\newenvironment{Exampleparagraph}[1][]
{
    \par 
    \noindent 
    \textbf{\color{blue} #1}
    \par 
    \vspace{0.5em} 
    \itshape 
}
{
    \par 
    \vspace{1em} 
}

%\input{local-macro}

\newcommand{\RED}[1]{\textcolor{red}{#1}}
\newcommand{\zc}[1]{\textcolor{red}{ZC: #1}}

\newcommand{\bern}[1]{\mathrm{Bern}\left(#1\right)}



\title{Discrete Proximal Sampler}
\date{\today}

\pgfplotsset{compat=1.18} 
\begin{document}

\maketitle

\section{Setup}
\begin{definition}[Discrete diffusion process]
Let $N_1(t),N_2(t),\ldots,N_n(t)$ be independent one-dimensional standard Poisson point processes.
Let $\mu$ be a distribution supported over $\{-1,+1\}^n$.
Starting from $X_0 \sim \mu$, the discrete diffusion process $(X_t)_{t \in [0,1]}$ is given by
\begin{align*}
    \forall i \in [n],\quad X_t(i) = 
    \begin{cases}
    X_t(i) & \text{if $N_i(-\frac{1}{2} \log (1-t))$ is even,}\\
    -X_t(i) & \text{otherwise.}
    \end{cases}
\end{align*}
The denoising process $(Y_t)_{t \in [0,1]} \overset{d}{=} (X_{1-t})_{t \in [0,1]}$ is the time-reversal of the discrete diffusion process.
\end{definition}

\begin{remark}
    The discrete diffusion process is a Markov process as $N_i$ is ``memoryless'' (in the sense that the distribution of $N_i(t+\Delta) - N_i(t)$ is independent of events before time $t$).
\end{remark}

\begin{definition}[Discrete proximal sampler]
Let $P_{\eta \to \theta}(\cdot,\cdot)$ and $Q_{\eta \to \theta}(\cdot,\cdot)$ be Markov kernels defined as follows:
\begin{align*}
    P_{\theta \to \eta}(x,y) = \Pr[]{X_{\eta} = y \mid X_{\theta} = x} \quad \text{and} \quad Q_{\eta \to \theta}(x,y) = \Pr[]{X_{1-\theta} = y \mid X_{1-\eta} = x}.
\end{align*}
The discrete proximal sampler is the Markov chain $P_{1 \to \eta} Q_{\eta \to 1}$.
\end{definition}

\begin{remark}
    Note that for each $i \in [n]$, the probability of $X_t(i) = X_0(i)$ is exactly $\frac{2-t}{2}$. Therefore,
    The posterior distribution $Q_{\eta \to 1}(x,\cdot)$ satisfies
    \begin{align*}
        Q_{\eta \to 1}(x,y) = \Pr[]{X_0 = y \mid X_{1-\eta} = x} &\propto \Pr[]{X_0 = y \text{ and } X_{1-\eta} = x}\\ &\propto \mu(y) \cdot \tp{\frac{1-\eta}{1+\eta}}^{\abs{x \oplus y }}.
    \end{align*}
\end{remark}

\section{Spectral Stability}
\begin{definition}[Spectral stability]
    The distribution $\mu$ satisfies the $C$-spectral stability with respect to the discrete diffusion process if for any $\eta \in [0,1)$, it holds
    \begin{align}\label{eq:def-SI}
        \forall x \in \+X \text{ and } f \in \mathbb{R}^{\+X},\quad \lim_{h \to 0^+} \frac{1}{h} \Var[Q_{\eta \to \eta + h}(x,\cdot)]{f_{\eta + h}} \le C(\eta) \cdot \Var[Q_{\eta \to 1}(x,\cdot)]{f},
    \end{align}
    where $f_\eta(y) = \E[Q_{\eta \to 1}(y,\cdot)]{f}$ and $\+X = \{-1,+1\}^n$ denotes the space of configurations.
\end{definition}

\begin{lemma}
    The condition~\eqref{eq:def-SI} is equivalent to
    \begin{align*}
        \frac{1}{2\eta} \cdot \sum_{k \in [n]} \frac{\Pr[]{X_{1-\eta} = x^{\oplus k}} }{\Pr[]{X_{1-\eta} = x}} \cdot \tp{\E[Q_{\eta \to 1}(x,\cdot)]{f} - \E[Q_{\eta \to 1}(x^{\oplus k},\cdot)]{f}}^2 \le C(\eta) \cdot \Var[Q_{\eta \to 1}(x,\cdot)]{f}.
    \end{align*}
    Specifically, when $\eta = 0$, the spectral stability holds with rate $C(0) = 0$.
\end{lemma}
\begin{proof}
The proof follows from a brute-force calculation. By definition, it holds
\begin{align*}
    Q_{\eta \to \eta+h}(x,y) = \Pr[]{X_{1-\eta-h} = y \mid X_{1-\eta}=x}&= \frac{\Pr[]{X_{1-\eta} = x \mid X_{1-\eta-h} = y} \cdot \Pr[]{X_{1-\eta-h} = y}}{\Pr[]{X_{1-
    \eta} = x}}\\
    &=\begin{cases}
    1+\Theta(h) & \text{if $x = y$}\\
    \frac{1}{2\eta} \cdot \frac{\Pr[]{X_{1-\eta} = y}}{\Pr[]{X_{1-\eta} = x}} \cdot h + o(h) & \text{$\abs{x \oplus y} = 1$}\\
    o(h) & \text{otherwise}.
    \end{cases}
\end{align*}
The last equation follows from that $\Pr[]{X_{1-\eta} = x \mid X_{1-\eta-h} = y}$ equals to the probability the difference between $N_{-\frac{1}{2} \log \eta}(k)$ and $N_{-\frac{1}{2} \log (\eta+h)}(k)$ is odd for sites $k \in x \oplus y$, and the difference is even for all $i \notin \abs{x \oplus y}$. Hence,
\begin{align*}
\Pr[]{X_{1-\eta} = x \mid X_{1-\eta-h} = y} = \tp{\frac{1+\frac{\eta}{\eta+h}}{2}}^{n-\abs{x \oplus y}} \cdot \tp{\frac{1-\frac{\eta}{\eta+h}}{2}}^{\abs{x \oplus y}}.
\end{align*}
Therefore, the variance $\Var[Q_{\eta \to \eta + h(x,\cdot)}]{f_{\eta+h}}$ satisfies
\begin{align*}
    \Var[Q_{\eta \to \eta+h}(x,\cdot)]{f_{\eta+h}} &= \sum_{y \in \{-1,+1\}^n} Q_{\eta \to \eta+h}(x,y) \cdot \tp{f_{\eta+h}(y) - \E[Q_{\eta \to 1}(x,\cdot)]{f}}^2\\
    &= \frac{h}{2\eta} \cdot \sum_{k \in [n]} \frac{\Pr[]{X_{1-\eta} = x^{\oplus k}} }{\Pr[]{X_{1-\eta} = x}} \cdot \tp{\E[Q_{\eta \to 1}(x,\cdot)]{f} - \E[Q_{\eta \to 1}(x^{\oplus k},\cdot)]{f}}^2 + o(h),
\end{align*}
where the last inequality follows from the following facts:
\begin{enumerate}
\item $Q_{\eta \to \eta+h}(x,y) = o(h)$ for all $y \in \{-1,+1\}^n$ with $\abs{x \oplus y} > 1$;
\item $f_{\eta+h}(x) - f_\eta(x) = f_{\eta+h}(x) - \E[Q_{\eta \to 1}(x,\cdot)]{f} = \Theta(h)$.
\end{enumerate}
\end{proof}

\section{Trickle-down Equation}
In this section, we establish the trickle-down equation for the discrete diffusion process.
Consider the total law of covariance for the conditional covariance $\mathrm{Cov}(Y_1 \mid Y_{\eta})$:
    \begin{align*}
        \mathrm{Cov}(Y_1 \mid Y_\eta) = \E[]{\mathrm{Cov}(Y_1 \mid Y_{\eta+h}) \mid Y_\eta} + \mathrm{Cov}\tp{\E[]{Y_1 \mid Y_{\eta + h}} \mid Y_\eta}.
    \end{align*}
    It now suffices to compute the covariance of $\E[]{Y_1 \mid Y_{\eta + h}}$ conditioned on $Y_\eta$. By definition,
    \begin{align*}
        &\mathrm{Cov}\tp{\E[]{Y_1 \mid Y_{\eta + h}} \mid Y_\eta = x}\\
        = &\sum_{z \in \{-1,+1\}^n} \Pr[]{Y_{\eta + h} = z \mid Y_{\eta} = x} \cdot \tp{\E[]{Y_1 \mid Y_{\eta + h} = z} - \E[]{Y_1 \mid Y_\eta = x}}^{\oplus 2}\\
        = & \frac{1}{2\eta} \cdot \sum_{k \in [n]} \frac{\Pr[]{X_{1-\eta} = x^{\oplus k}} }{\Pr[]{X_{1-\eta} = x}} \cdot \tp{\*m(Q_{\eta \to 1}(x^{\oplus k},\cdot)) - \*m(Q_{\eta \to 1}(x,\cdot))}^{\otimes 2} \cdot h+ o(h)
    \end{align*}
    It turns out that the linear term is closely related to the covariance matrix.
    \begin{proposition}
        For any $x \in \{-1,+1\}^n$, it holds
        \begin{align}\label{eq:target-trickle-down}
        \frac{\Pr[]{X_{1-\eta} = x^{\oplus k}}}{\Pr[]{X_{1-\eta} = x}} \cdot \tp{\*m_i(Q_{\eta \to 1}(x^{\oplus k},\cdot)) - \*m_i(Q_{\eta \to 1}(x,\cdot))} = \frac{4\eta \cdot\mathrm{sign}(x_ix_k)}{(1-\eta)(1+\eta)} \cdot \mathrm{Cov}(Q_{\eta \to 1}(x,\cdot))_{i,k}.
        \end{align}
    \end{proposition}

    \begin{proof}
        Fix $x \in \{-1,+1\}^n$. Let $\nu$ denote the distribution $Q_{\eta \to 1}(x,\cdot)$, and $R = \frac{\Pr[]{X_{1-\eta} = x^{\oplus k}}}{\Pr[]{X_{1-\eta} = x}}$. By the law of total probability, it holds
        \begin{align*}
            \Pr[]{X_{1-\eta} = x^{\oplus k}} &= \Pr[]{X_{1-\eta} = x^{\oplus k} \text{ and } X_0(k) = x_k} + \Pr[]{X_{1-\eta} = x^{\oplus k} \text{ and } X_0(k) = -x_k} \\
            &= \frac{1-\eta}{1+\eta} \cdot \Pr[]{X_{1-\eta} = x \text{ and } X_0(k) = x_k} + \frac{1+\eta}{1-\eta} \cdot \Pr[]{X_{1-\eta} = x \text{ and } X_0(k) = -x_k} \\
            &=\Pr[]{X_{1-\eta} = x} \cdot \tp{\frac{1+\eta}{1-\eta} - \frac{4\eta}{(1-\eta)(1+\eta)} \cdot \nu\tp{X_k = x_k}}.
        \end{align*}
        Similarly, we have 
        \begin{align*}
            &\Pr[]{X_0(i) = x_i \text{ and } X_{1-\eta} = x^{\oplus k}} \\
            = &\Pr[]{X_{1-\eta} = x} \cdot \tp{\frac{1-\eta}{1+\eta} \cdot \nu(X_i = x_i \text{ and } X_k = x_k) - \frac{1+\eta}{1-\eta} \cdot \nu(X_i = x_i \text{ and } X_k=-x_k)}.
        \end{align*}
        The left-hand-side of~\eqref{eq:target-trickle-down} equals to
        \begin{align*}
            \text{LHS of~\eqref{eq:target-trickle-down}} &= \tp{R-\frac{1-\eta}{1+\eta}} \cdot \nu\tp{X_i=x_i \text{ and } X_k = x_k}\\ &+ \tp{R - \frac{1+\eta}{1-\eta}} \cdot \nu\tp{X_i = x_i \text{ and } X_k=-x_k}\\
            &=\frac{4\eta}{(1+\eta)(1-\eta)}\tp{\nu\tp{X_i = x_i \text{ and } X_k=x_k} - \nu(X_i = x_i)\cdot \nu(X_k=x_k)}\\
            &=\frac{4\eta \cdot x_i x_k}{(1-\eta)(1+\eta)} \cdot \mathrm{Cov}(\nu)_{i,k}.
        \end{align*}
    \end{proof}
    As a straightforward corollary, we have the following trickle-down equation.
    \begin{lemma}[Trickle-down equation for discrete diffusion process]
    Fix $\eta \in [0,1)$ and let $\Sigma$ denote the covariance matrix $\mathrm{Cov}(Y_1 \mid Y_h=x)$. It holds
    \begin{align*}
        \Sigma = \E[]{\mathrm{Cov}(Y_1 \mid Y_{\eta+h}) \mid Y_\eta=x} + \frac{8\eta}{(1-\eta^2)^2} \cdot \mathrm{diag}\tp{x} \cdot \Sigma \Pi_{\eta,x}^{-1} \Sigma \cdot \mathrm{diag} \tp{x} \cdot h + o(h),
    \end{align*}
    where $\Pi_{\eta,x}$ is the diagonal matrix $ \mathrm{diag}\tp{\frac{\Pr[]{X_{1-\eta} = x^{\oplus k}}}{\Pr[]{X_{1-\eta} = x}}}_{1 \le k \le n}$.
    \end{lemma}

    To establish the trickle-down theorem for discrete diffusion process, we prove the following fact. 
    \begin{proposition}
        Let $0 \le \eta \le \theta \le 1$ and $x \in \{-1,+1\}^n$. It holds
        \begin{align}\label{eq:expectation}
            \E[]{\Pi_{\theta,Y_{\theta}} \mid Y_\eta = x} = \Pi_{\eta,x}.
        \end{align}
    \end{proposition}
    \begin{proof}
         For each $k \in [n]$, it holds
        \begin{align*}
            &\sum_{y \in \{-1,+1\}^n} \Pr[]{Y_{\theta} = y \mid Y_\eta = x} \cdot \frac{\Pr[]{Y_{\theta} = y^{\oplus k}}}{\Pr[]{Y_\theta = y}}\\
            =& \frac{1}{\Pr[]{Y_\eta =x}} \cdot \sum_{y \in \{-1,+1\}^n} \Pr[]{Y_\eta = x \mid Y_\theta = y} \cdot \Pr[]{Y_\theta = y^{\oplus k}}\\
            =&\frac{1}{\Pr[]{Y_\eta = x}} \cdot \sum_{y \in \{-1,+1\}^n} \Pr[]{Y_\eta = x^{\oplus k} \mid Y_\theta = y^{\oplus k}} \cdot \Pr[]{Y_\theta = y^{\oplus k}}\\
            =& \frac{\Pr[]{Y_\eta = x^{\oplus k}}}{\Pr[]{Y_\theta = x}}.
        \end{align*}
        This implies~\eqref{eq:expectation}
    \end{proof}

\section{Trickle-down Theorem}
Suppose $C:[0,1]\to \mathbb{R}$ is a smooth function such that $\mathrm{Cov}(Y_1 \mid Y_\eta) \preceq C(\eta) \cdot \Pi_{\eta,Y_\eta}$. The trickle-down equation suggests the following differential equation on $C$:
\begin{equation}\label{eq:C_2spin}
    C(\eta) \le C(\eta+h) + \frac{8\eta}{(1-\eta^2)^2} C(\eta)^2 \cdot h +o(h).
\end{equation}
The solution to this differential equation is $C(\eta) \le U(\eta) := \frac{1}{K+\frac{4}{1-\eta^2}}$, where $K > -4$ is a constant determined by boundary condition. Note that
\begin{align*}
    U(1-h) = \frac{1}{2}\cdot h - \frac{1+K}{4}\cdot h^2 + o(h^2).
\end{align*}
It suggests us to investigate the covariance matrix $\mathrm{Cov}(Y_1 \mid Y_{1-h} = y)$ and diagonal matrix $\Pi_{1-h,y}$ up to the second order.

\begin{proposition}
For all $y \in \{-1,+1\}^n$ and $i \in [n]$, it holds
\begin{align*}
    \Pr[X \sim Q_{(1-h) \to 1}(y,\cdot)]{X_i \neq y_i} = \frac{\mu(y^{\oplus i})}{2 \mu(y)} \cdot h + \frac{1}{4} \cdot \tp{\frac{\mu(y^{\oplus i})}{\mu(y)} + \sum_{k \neq i} \frac{\mu(y^{\oplus \{k,i\}})}{\mu(y)} - \frac{\mu(y^{\oplus i})\sum_j\mu(y^{\oplus j})}{\mu(y)^2}} \cdot h^2 + o(h^2).
\end{align*}
\end{proposition}

\begin{proof}
    Note that the probability of $X \sim Q_{(1-h) \to 1}(y,\cdot)$ is in the order of $o(h^2)$ if the symmetric difference $\abs{X \oplus y}$ is greater than $2$. Therefore, it holds
    \begin{align*}
    \Pr[X \sim Q_{(1-h) \to 1}(y,\cdot)]{X_i \neq y_i} &= \frac{\frac{h}{2-h}\mu(y^{\oplus i}) + \frac{h^2}{(2-h)^2} \sum_{k \neq i} \mu(y^{\oplus \{k,i\} }) }{\mu(y) + \frac{h}{2-h}\sum_{j} \mu(y^{\oplus j}) + \frac{h^2}{(2-h)^2} \sum_{1 \le j < k \le n} \mu(y^{\oplus \{j,k\}})} + o(h^2)\\
    &= \frac{\mu(y^{\oplus i})}{2 \mu(y)} \cdot h + \underbrace{\frac{1}{4} \cdot \tp{\frac{\mu(y^{\oplus i})}{\mu(y)} + \sum_{k \neq i} \frac{\mu(y^{\oplus \{k,i\}})}{\mu(y)} - \frac{\mu(y^{\oplus i}) \sum_{j} \mu(y^{\oplus j}) }{\mu(y)^2}}}_{:=S_i} \cdot h^2 + o(h^2)
\end{align*}
\end{proof}
We are now ready to calculate the covariance matrix and the diagonal matrix.
\begin{align*}
    \Pi_{1-h,x}(i,i) &= \frac{h}{2-h} \cdot \Pr[X \sim Q_{(1-h) \to 1}(y,\cdot)]{X_i = y_i} + \frac{2-h}{h} \cdot \Pr[X \sim Q_{(1-h) \to 1}(y,\cdot)]{X_i \neq y_i}\\
    &=\frac{h}{2} + \tp{\frac{2-h}{h} - \frac{h}{2-h}} \cdot \Pr[X \sim Q_{(1-h) \to 1}(y,\cdot)]{X_i \neq y_i} + o(h)\\
    &=\frac{h}{2} + \frac{1}{h} \cdot \tp{2-h+o(h)} \cdot \tp{\frac{\mu(y^{\oplus i})}{2\mu(y)} \cdot h + S_i h^2 +o(h^2)} +o(h)\\
    &= \frac{\mu(y^{\oplus i})}{\mu(y)} + \tp{\frac{1}{2} + 2S_i - \frac{\mu(y^{\oplus i})}{2\mu(y)}} \cdot h + o(h).
\end{align*}
Similarly, it holds
\begin{align*}
    \mathrm{Cov}(Y_1 \mid Y_{1-h}=y)(i,j) = (\frac{\mu(y^{\oplus \{i,j\}})}{4\mu(y)} - \frac{\mu(y^{\oplus i})\mu(y^{\oplus j})}{4\mu(y)^2}) \cdot h^2+o(h^2)
\end{align*}
and 
\begin{align*}
\mathrm{Cov}(Y_1 \mid Y_{1-h}=y)(i,i) = \frac{\mu(y^{\oplus i})}{2 \mu(y)} \cdot h + \tp{S_i - \frac{\mu(y^{\oplus i})^2}{4\mu(y)^2}} \cdot h^2 +o(h^2).
\end{align*}
To obtain the trickle-down theorem, it suffices to show that
\begin{align*}
    \mathrm{Cov}\tp{Y_1 \mid Y_{1-h}=y} -\frac{h}{2}\cdot \Pi \preceq \frac{3}{4} \cdot \Pi \cdot h^2 +o(h^2),
\end{align*}
which is equivalent to that
\begin{align*}
    M \preceq 2 \cdot \mathrm{diag}\tp{\frac{\mu(y^{\oplus i})}{\mu(y)}},
\end{align*}
where $M$ is the matrix such that
\begin{align*}
    M(i,j) = \frac{\mu(y^{\oplus \{i,j\}})}{\mu(y)} - \frac{\mu(y^{\oplus i})\mu(y^{\oplus j})}{\mu(y)^2} \quad \text{and} \quad M(i,i) = - 1 - \frac{\mu(y^{\oplus i})^2}{\mu(y)^2}.
\end{align*}


\zc{
Using similar notations as in [CCCYZ25]: Let
\begin{align*}
    M(i,j) := 
    \begin{cases}
        \frac{\mu(y^{\oplus \{i,j\}}) \mu(y)}{\mu(y^{\oplus i}) \mu(y^{\oplus j})} - 1, & i \neq j;\\
        0, & i = j.
    \end{cases}
\end{align*}
Then the condition is
\begin{align*}
    M \preceq \mathrm{diag}\tp{ \left( 1+\frac{\mu(y)}{\mu(y^{\oplus i})} \right)^2 }
\end{align*}
}

\section{$q$-spin System}
\subsection{Similar Setup}
In the same way, we can define the discrete diffusion process.
The only the difference is that every time the Poisson clock rings, we pick a color $c \in [q]$ uniformly at random.
We have that
$$
    Q_{\eta \to 1}(x,y) \propto \left(\frac{1-\eta}{q\eta + 1 - \eta}\right)^{d_H(x,y)} \mu(y);
$$
$$
    \Pr{Y_\eta = X \mid Y_{\eta + h} = Y}
    =\left(\frac{h+q\eta}{q(\eta+h)}\right)^{n} \left(\frac{h}{q\eta+h}\right)^{d_H(X,Y)} .
$$
\subsection{Spectral Stability}
\begin{lemma}\label{lem:qspin_SS1}
    The condition~\eqref{eq:def-SI} is equivalent to
    \begin{align*}
        \frac{1}{q\eta} \sum_{i,c\neq x_i} \frac{\Pr{Y_{\eta} = x^{i\to c}}}{\Pr{Y_\eta = x}} \left(f_{\eta}(x^{i\to c}) - \E[Q_{\eta\to1}(x,\cdot)]{f}\right)^2 \leq C(\eta) \cdot \Var[Q_{\eta \to 1}(x,\cdot)]{f}.
    \end{align*}
\end{lemma}
\begin{proof}
    We have that
    $$
        Q_{\eta \to \eta + h}(x,y) = \begin{cases}
            \Omega(1) , & x = y\\
            \frac{h}{q\eta}\cdot\frac{\Pr{Y_{\eta + h } = y}}{\Pr{Y_\eta = x}} + o(h) , & d_H(x,y) = 1\\
            o(h), & \text{otherwise}
        \end{cases} .
    $$
    And for any $x\in [q]^V$, 
    $$
        f_{\eta+h}(x) - f_\eta(x) = \sum_{y\in [q]^V} f(y)(Q_{\eta+h\to1}(x,y) - Q_{\eta\to1}(x,y)) = O(h).
    $$
    Therefore, the variance $\Var[Q_{\eta\to\eta+h}(x,\cdot)]{f_{\eta+h}}$ is as follows.
    \begin{align*}
        \Var[Q_{\eta\to\eta+h}(x,\cdot)]{f_{\eta+h}} &= \sum_{i,c\neq x_i} Q_{\eta\to\eta+h}(x,y) \left(f_{\eta+h}(x^{i\to c}) - f_\eta(x)\right)^2 + o(h)
        \\&= \frac{h}{q\eta} \sum_{i,c\neq x_i} \frac{\Pr{Y_{\eta} = x^{i\to c}}}{\Pr{Y_\eta = x}} \left(f_{\eta+h}(x^{i\to c}) - f_\eta(x)\right)^2 + o(h) .
        \qedhere
    \end{align*}
\end{proof}
Let $\frac{\nu}{Q_{\eta\to 1}(x,\cdot)} := \frac{f}{f_\eta(x)}$ be the relative density.
We have that 
$\Var[Q_{\eta\to 1}(x,\cdot)]{f} = D_{\chi^2}(\nu\parallel\mu)f_\eta(x)^2$.
Let $r_{i,c} = \frac{\Pr{Y_\eta = x^{i\to c}}}{\Pr{Y_\eta = x}}$.
The following holds:
\begin{align*}
    \frac{Q_{\eta\to 1}(x^{i\to c},y)}{Q_{\eta\to 1}(x,y)} = 
    \frac{\Pr{Y_\eta = x^{i\to c}\mid Y_1 = y}\Pr{Y_\eta = x}}{\Pr{Y_\eta = x^{}\mid Y_1 = y}\Pr{Y_\eta = x^{i\to c}}}
    = r_{i,c}^{-1} \frac{dif(c,y_i)}{dif(x_i,y_i)}.
\end{align*}
Therefore,
\begin{align}
    f_\eta(x^{i\to c}) &= \sum_{T}Q_{\eta\to 1}(x^{i\to c},T) f(T)
    = r^{-1}_{i,c}\sum_{T}Q_{\eta\to 1}(x,T) f(T) \frac{dif(c,T_i)}{dif(x_i,T_i)} \nonumber
    \\&= r_{i,c}^{-1}f_\eta(x)\sum_{c_2\in [q]}\frac{dif(c,c_2)}{dif(x_i,c_2)} \*m(\nu)_{i,c_2} \label{eq:qspin_SS}
\end{align}
Plugging Lemma \ref{lem:qspin_SS1} and Equation \ref{eq:qspin_SS}, we have the following proposition.
\begin{proposition}
    Let $\gamma:= Q_{\eta\to 1}(x,\cdot)$. The condition~\eqref{eq:def-SI} is equivalent to 
    $$
        \forall \nu \ll \gamma: \qquad \frac{1}{q\eta}\sum_{i,c\neq x_i}r_{i,c} (r_{i,c}^{-1}\sum_{c_2\in [q]}\frac{dif(c,c_2)}{dif(x_i,c_2)} \*m(\nu)_{i,c_2} - 1)^2 \leq C(\eta)\cdot D_{\chi^2}(\nu \pll \gamma) .
    $$
\end{proposition}
\subsection{Trickle-down Equation}
In the same way, we only need to simplify $\Cov(\E{Y_1 \mid Y_{\eta+ h}}\mid Y_\eta = x)$.
By definition,
\begin{align}
    \Cov(\E{Y_1\mid Y_{\eta+ h}}\mid Y_\eta = x)
    &= \sum_{z\in [q]^V} \Pr{Y_{\eta+h} = z\mid Y_\eta = x} \cdot 
    (\E{Y_1\mid Y_{\eta + h} = z} - \E{Y_1\mid Y_\eta = x})^{\otimes 2}. \nonumber
    \\&= \sum_{i,c\neq x_i} \frac h{q\eta} \frac{\Pr{Y_\eta = x^{i\to c}}}{\Pr{Y_\eta = x}}(\E{Y_1\mid Y_{\eta + h} = x^{i\to c}} - \E{Y_1\mid Y_\eta = x})^{\otimes 2} + o(h). \nonumber
    \\&= \frac{h}{q\eta}\sum_{i,c\neq x_i} \frac{\Pr{Y_\eta = x^{i\to c}}}{\Pr{Y_\eta = x}} \left(\*m(Q_{\eta\to 1}(x^{i\to c},\cdot)) - \*m(Q_{\eta\to 1}(x,\cdot))\right)^{\otimes 2} + o(h) \label{eq:trickle_down_base}
\end{align}
We use $r_{i,c}$ to denote $\frac{\Pr{Y_\eta = x^{i\to c}}}{\Pr{Y_\eta = x}}$ for simplicity.

Let $\nu = Q_{\eta \to 1}(x,\cdot)$ and $\nu' = Q_{\eta \to 1}(x^{i\to c},\cdot)$.
Let $dif: q\times q \mapsto \R$ and $dif(x,y) = \frac 1q\begin{cases}1-\eta,&x\neq y\\q\eta+1-\eta, &x=y\end{cases}$.

\begin{proposition}\label{prop:q_spin_Cov}
    Given $j\in[n]$ and $\overline{c}\in [q]$, we have that 
    $$
        r_{i,c} (m_{j,\overline c}(\nu') - m_{j,\overline c}(\nu) ) = \sum_{c' \in [q]} \frac{dif(c,c')}{dif(x_i,c')} \Cov_\nu((j,\overline c),(i,c')).
    $$
\end{proposition}
\begin{proof}
    We have that for any $i\in[n]$ and $c\in[q]$,
    \begin{align}
        \nu'(y) &= \Pr{Y_1 = y\mid Y_\eta = x^{i\to c}}
        = \Pr{Y_\eta = x^{i\to c}\mid Y_1 = y}\frac{\Pr{Y_1 = y}}{\Pr{Y_\eta = x^{i\to c}}}\nonumber 
        \\&= \frac{dif(y_i,c)}{dif(y_i,x_i)}\Pr{Y_\eta = x\mid Y_1 = y}\frac{\Pr{Y_1 = y}}{r_{i,c}\Pr{Y_\eta =x}}
        = \frac{dif(y_i,c)}{dif(y_i,x_i)r_{i,c}} \nu(y) \label{eq:q_spin_nu'}.
    \end{align}
    By \Cref{eq:q_spin_nu'}, 
    \begin{equation}\label{eq:q_spin_eq_1}
        m_{j,\overline c}(\nu') = \sum_{c'\in[q]}\frac{dif(c',c)}{dif(c',x_i)r_{i,c}}\nu(\sigma_j=\overline c,\sigma_i = c'). 
    \end{equation}
    By law of total probability, it holds that 
    \begin{align*}
        \Pr{Y_\eta = x^{i\to c}} 
        &= \sum_{c'\in[q]} \Pr{Y_\eta = x^{i\to c}, Y_1(i) = c'} \nonumber
        \\&= \Pr{Y_\eta = x^{i\to c},Y_1(i) = x_i} + \Pr{Y_\eta = x^{i\to c},Y_1(i) = c} + \sum_{c'\neq x_i\land c'\neq c} \Pr{Y_\eta = x, Y_1(i) = c'} \nonumber
        \\&= \Pr{Y_\eta = x}(\frac{dif(c,x_i)}{dif(x_i,x_i)}\nu(\sigma_i = x_i)+ \frac{dif(c,c)}{dif(x_i,c)}\nu(\sigma_i = c) + \sum_{c'\neq x_i\land c'\neq c} \nu(\sigma_i = c') )  ,
    \end{align*}
    which implies that 
    \begin{equation}\label{eq:q_spin_eq_2}
        r_{i,c} = \sum_{c'\in [q]} \frac{dif(c,c')}{dif(x_i,c')} \nu(\sigma_i = c')  .
    \end{equation}
    Plugging in \Cref{eq:q_spin_eq_1} and \Cref{eq:q_spin_eq_2}, we have that 
    \[
        r_{i,c} (m_{j,\overline c}(\nu') - m_{j,\overline c}(\nu) ) = \sum_{c' \in [q]} \frac{dif(c,c')}{dif(x_i,c')} (\nu(\sigma_j=\overline c,\sigma_i = c') - \nu(\sigma_j=\overline c)\nu(\sigma_i = c')) .
        \qedhere 
    \]
\end{proof}

\begin{lemma}[Trickle-down Equation on $q$-spin systems]\label{lem:trickle_down_q}
    Let $\*d_c = (\frac{dif(c,1)}{dif(x_i,1)},\frac{dif(c,2)}{dif(x_i,2)},\ldots,\frac{dif(c,q)}{dif(x_i,q)})^\top$.
    Then we have that
    \begin{align*}
        \mathrm{Cov}(Y_1 \mid Y_\eta = x) = \E[]{\mathrm{Cov}(Y_1 \mid Y_{\eta+h}) \mid Y_\eta = x} + \frac{h}{q\eta} \Cov(Y_1\mid Y_\eta = x)\*B\Cov(Y_1\mid Y_\eta = x) + o(h),
    \end{align*}
    where
    $$
        \*B = \diag\{B_i = \sum_{c\neq x_i} \frac{\Pr{Y_\eta = x}}{\Pr{Y_\eta = x^{i\to c}}}\*d_c\*d_c^\top\}_{i\in[n]}.%\quad \text{and} \quad \Pi_x = \diag\{dif(x_i,c)\}_{(i,c)\in [n]\times [q]}.
    $$
\end{lemma}
\begin{proof}
    By \Cref{eq:trickle_down_base} and Propsition \ref{prop:q_spin_Cov},
    \begin{align*}
        &\Cov(\E{Y_1\mid Y_{\eta+ h}}\mid Y_\eta = x)((j_1,c_1),(j_2,c_2))
        \\=&\frac{h}{q\eta}\sum_{i,c\neq x_i} \frac 1{r_{i,c}} (\sum_{c' \in [q]} \frac{dif(c,c')}{dif(x_i,c')} \Cov_\nu((j_1,c_1),(i,c'))) (\sum_{c' \in [q]} \frac{dif(c,c')}{dif(x_i,c')} \Cov_\nu((i,c'),(j_2,c_2))) + o(h)
        \\=&\frac{h}{q\eta}\sum_{i,c_1',c_2'}\frac{\Cov_\nu((j_1,c_1),(i,c'_1))\Cov_\nu((i,c'_2),(j_2,c_2))}{dif(x_i,c_1')dif(x_i,c_2')} \sum_{c\neq x_i} \frac {dif(c,c_1')dif(c,c_2')}{r_{i,c}} + o(h) .
    \end{align*}
    Since $\Cov(Y_1\mid Y_h = x) = \Cov(\nu)$, we conclude the proof by the law of total covariance.
\end{proof}
% \begin{proposition}
%     Fix $i\in[n]$, let $v = (\frac{\Pr{Y_\eta = x}}{\Pr{Y_\eta = x^{i\to c}}} - \ind[c=x_i])_{c\in[q]}$, $k_1 = 1 - \frac{1-\eta}{q-1+\eta}$ and $k_2 = \frac{q-1+\eta}{1-\eta} -1$. We have that
%     $$
%         \Pi_{i,x}^{-1}B_i\Pi_{i,x}^{-1} = k_2^2 \diag\{v\} + \begin{pmatrix}1 &k_2 v &k_1 e_{x_i}\end{pmatrix}
%         \begin{pmatrix}
%             2s &1 &-s
%             \\1 & 0 & 1
%             \\-s & 1 & 2s
%         \end{pmatrix}
%         \begin{pmatrix}1 &k_2 v &k_1 e_{x_i}\end{pmatrix}^\top ,
%     $$
%     where $s = \sum_{c\neq x_i} \frac{\Pr{Y_\eta = x}}{\Pr{Y_\eta = x^{i\to c}}}$ .
%     The eigenvalues of the $3\times 3$ matrix are $3s, \frac{s\pm\sqrt{s^2+8}}{2}$.
% \end{proposition}

\begin{conjecture}\label{conj:induction}
    The induction inequality (Spectral Stability) of trickle-down is as follows:
    $$
        \Cov(Y_1\mid Y_\eta = x) \preceq C(\eta)\cdot \diag\{\sum_{c\neq x_i}\frac{\Pr{Y_\eta = x^{i\to c}}}{\Pr{Y_\eta = x}} a_c a_c^\top  \}_{i\in[n]}, 
    $$
    where $a_c\in \-{span}\{d_1,\ldots d_{c-1},d_{c+1},\ldots,d_q,\*1\}^{\perp}$.
    Formally, we can assume that 
    $$
        a_c(r) = \begin{cases}
            \alpha + (q-2)\beta &, r = c
            \\-\alpha &, r = x_i
            \\-\beta &, o.w.
        \end{cases}\quad \text{and}\quad a_c^\top d_c = \alpha(\alpha + (q-1)\beta) .
    $$
    where $\alpha = \frac{q\eta}{1-\eta}$, $\beta = \frac{q\eta}{q\eta + 1-\eta}$ and $d_c = 1 + \alpha e_c - \beta e_{x_i}$.
    % Specifically, $q=2$ implies $a_c^\top d_c = \eta$.
\end{conjecture}
If $q = 2$ and $x = (0,0)^\top$, $a_c = \alpha(-1,1)^\top$.
Plugging Lemma \ref{lem:trickle_down_q} and Conjecture \ref{conj:induction}, let $\tilde C(\eta) = C(\eta)\alpha^2$, the induction inequality is as follows,
$$
    C(\eta) \leq \tilde C(\eta + h) + \frac{h}{q\eta}(a_c^\top d_c)^2 \tilde C(\eta)^2 \frac{1}{\alpha^2} = \tilde C(\eta + h) + \frac{8\eta h}{(1-\eta^2)^2}\tilde C(\eta)^2,
$$ 
which is exactly Equation \ref{eq:C_2spin}.


Plugging Conjecture \ref{conj:induction}, we have that 
\begin{align}
    &\E{\Cov(Y_1\mid Y_{\eta+h})\mid Y_\eta=x} \nonumber
    \\&= \sum_{y\in[q]^V} \Cov(Y_1\mid Y_{\eta+h}=y)\Pr{Y_{\eta+h}=y\mid Y_\eta =x} \nonumber
    \\&\leq C(\eta+h)\sum_y \diag\{\sum_{c\neq y_i}\frac{\Pr{Y_{\eta+h} = y^{i\to c}}}{\Pr{Y_{\eta+h} = y}} a_c a_c^\top  \}_{i\in[n]}Q_{\eta\to\eta + h}(x,y) \nonumber
    % \\&=C(\eta+h)\sum_{j,c\neq x_j} \diag\{\sum_{c_2\neq x_i^{j\to c}}\frac{\Pr{Y_{\eta+h} = x^{i\to c_2,j\to c}}}{\Pr{Y_{\eta} = x}} a_c a_c^\top  \}_{i\in[n]}\cdot h
    % \\&+C(\eta+h)\diag\{\sum_{c_2\neq x_i}\frac{\Pr{Y_{\eta+h} = x^{i\to c_2}}}{\Pr{Y_{\eta+h} = x}} a_c a_c^\top  \}_{i\in[n]}\cdot (1 - \sum_{j,c\neq x_j}\frac{\Pr{Y_{\eta+h}=x^{j\to c}}}{\Pr{Y_{\eta}=x}}h)
    \\&= \frac{C(\eta+h)}{\Pr{Y_\eta=x}}\sum_y \diag\{\sum_{c\neq y_i}\Pr{Y_{\eta+h} = y^{i\to c}} a_{c,y,\eta+h} a_{c,y,\eta+h}^\top  \}_{i\in[n]}\Pr{Y_\eta =x\mid Y_{\eta+h}=y} \label{eq:1} .
\end{align}
Fix $i\in[n]$.
Since $\alpha(\eta + h) - \alpha(\eta) = \frac{qh}{(1-\eta)^2} + o(h) = \alpha(\eta)\frac{h}{\eta(1-\eta)} + o(h)$ and
$\beta(\eta + h) - \beta(\eta) = \frac{qh}{(q\eta+1-\eta)^2} + o(h) = \beta(\eta)\frac{h}{\eta(q\eta+1-\eta)} + o(h)$,
Equation \ref{eq:1} is equivalent to
\begin{align*}
    &\sum_y \sum_{c\neq y_i}\Pr{Y_{\eta+h} = y^{i\to c}} a_{c,y_i,\eta+h} a_{c,y_i,\eta+h}^\top \Pr{Y_\eta =x\mid Y_{\eta+h}=y}
    \\=&\sum_y (\frac{h}{q\eta+h})^{\ind[y_i\neq x_i]}\sum_{c\neq y_i}\Pr{Y_{\eta+h} = y^{i\to c}} a_{c,y_i,\eta+h} a_{c,y_i,\eta+h}^\top \Pr{Y_\eta =x^{i\to c}\mid Y_{\eta+h}=y^{i\to c}}  
    \\=&\frac{h}{q\eta}\sum_{y_i\neq x_i,c\neq y_i} \Pr{Y_{\eta+h} = x^{i\to c}} a_{c,y_i,\eta} a^\top_{c,y_i,\eta} \Pr{Y_\eta =x^{i\to c}\mid Y_{\eta+h}=x^{i\to c}}
    \\+&\sum_{c\neq x_i} a_{c,x_i,\eta+h} a^\top_{c,x_i,\eta+h} \Pr{Y_\eta =x^{i\to c}, Y_{\eta+h}=x^{i\to c}}
    \\+&\sum_{j\neq i,c_2\neq x_j}\sum_{c\neq x_i} a_{c,x_i,\eta} a^\top_{c,x_i,\eta} \Pr{Y_\eta =x^{i\to c}, Y_{\eta+h}=x^{i\to c,j\to c_2}} + o(h)
    \\=&\sum_{c\neq x_i} a_{c,x_i,\eta+h} a^\top_{c,x_i,\eta+h} \Pr{Y_\eta =x^{i\to c}} + \frac{h}{q\eta}\sum_{y_i\neq x_i,c\neq y_i} \Pr{Y_{\eta} = x^{i\to c}} a_{c,y_i,\eta} a^\top_{c,y_i,\eta}
    \\-&\sum_{c\neq x_i}\sum_{c_2\neq c} a_{c,x_i,\eta} a^\top_{c,x_i,\eta} \Pr{Y_\eta =x^{i\to c}, Y_{\eta+h}=x^{i\to c_2}} + o(h) ,
\end{align*}
where $\Pr{Y_\eta =x^{i\to c}, Y_{\eta+h}=x^{i\to c_2}} = \frac{h}{q\eta}(1+o(1))\Pr{Y_{\eta}=x^{i\to c_2}}$.
For $q = 2$, we have that 
$$
    \forall c\in[q], \sum_{y_i\neq x_i,y_i\neq c} a_{c,y_i}a_{c,y_i}^\top - \sum_{c_2\neq x_i,c_2\neq c}a_{c_2,x_i}a_{c_2,x_i}^\top = 0.
$$
since $a_{1,2}\parallel a_{2,1}$.
Then the induction works.
For $q > 2$, take $c_1\neq c_3$, we have that 
$$
\forall c\in[q], d_{c_1,x_i}^\top\sum_{y_i\neq x_i,y_i\neq c} a_{c,y_i}a_{c,y_i}^\top d_{c_3,x_i} \neq 0 =  d_{c_1,x_i}^\top \sum_{c_2\neq x_i,c_2\neq c}a_{c_2,x_i}a_{c_2,x_i}^\top  d_{c_3,x_i}.
$$

\end{document}  